\SourceLeaf{314}%
in den Modulus der briggischen Logarithmen bezeichnen, \(\rho\) in der im Art. 16 angegebenen Bedeutung genommen (und damit auch \(\mu\)). Man hat für diese constanten Factoren
\[
\begin{aligned}
\log\mu&=7.9297527989\,(-20)\\
\log\nu&=4.9206912908\,(-30)
\end{aligned}
\]

\indent Bis zu den Gliedern zweiter Ordnung stimmen diese Resultate mit den im 17. Art. gegebenen überein. Der Zweck der vorstehenden weitern Entwicklung war nur, klar hervortreten zu lassen, dass selbst zur schärfsten Rechnung die Glieder zweiter Ordnung völlig zureichen: in der That kommt in dem ganzen Hannoverschen Dreieckssysteme kein Fall vor, wo die Glieder vierter Ordnung den Betrag von zwei Einheiten der zehnten Decimale erreichten, und nur ein Paar Fälle, wo sie \textit{Eine} Einheit der zehnten Decimale überschreiten.

\Article{20}
Wenn unsere Formeln. welche nicht von der Breite und dem Azimuth an dem einen Orte, sondern von dem Mittelwerthe dieser Grössen an den beiden Örtern ausgehen, zur Auflösung der zu Anfang dieser Abhandlung aufgeführten Aufgabe benutzt werden sollen, so wird diess auf eine indirecte Art, oder richtiger durch stufenweise beliebig weit getriebene Annäherung geschehen müssen. Der Gang der Arbeit besteht darin, dass man von irgend einem genäherten Werthe von \(T\) ausgeht (wofür man in Ermangelung aller anderweitigen Kenntniss oder Schätzung zuerst das gegebene Azimuth an dem ersten Orte annehmen kann) und daraus einen viel schärfern ableitet; mit diesem dann dieselbe Rechnung wiederholt, und damit so lange fortfährt, bis man zu stehenden Resultaten gelangt. Man hat dabei zu beachten, dass bei den ersten Rechnungen nur 4 oder 5 Zifern der Logarithmen berücksichtigt zu werden brauchen, und dabei \(\beta\) und \(\tau\) anstatt der corrigirten \(b\) und \(t\) angewandt werden dürfen, daher man auch, bei diesen ersten Rechnungen, sich um \(\lambda\) und \(l\) noch nicht zu bekümmern braucht. Die Formeln sind so, wenn für den ersten Ort die Breite mit \(B^0\), das Azimuth mit \(T^0\) bezeichnet wird, der Reihe nach folgende:
\[
\begin{aligned}
\beta&=r\cos T\\
B&=B^0-\tfrac12\beta\\
\tau&=r\sin T\operatorname{tang}B\\
T&=T^0-\tfrac12\tau
\end{aligned}
\]
